\subsection{The Grothendieck Group of an Additive Set of Modules}

Let A be a ring. In this subsection, we consider an additive set C of classes of A-modules; we identify C with the canonical basis of the free abelian group \(\mathbf{Z}^{(\mathcal{C})}\) . For any exact sequence

\[
0 \longrightarrow \mathrm{E} ^ {\prime} \longrightarrow \mathrm{E} \longrightarrow \mathrm{E} ^ {\prime \prime} \longrightarrow 0\tag{\( \mathcal{E} \)}
\]

of modules of type C, we denote by \(r_{\mathcal{E}}\) the element \(\mathrm{cl}(\mathrm{E}) - \mathrm{cl}(\mathrm{E}') - \mathrm{cl}(\mathrm{E}'')\) of \(\mathbf{Z}^{(\mathcal{C})}\) . Let R be the subgroup of \(\mathbf{Z}^{(\mathcal{C})}\) generated by the elements of the form \(r_{\mathcal{E}}\) ; the quotient group \(\mathbf{Z}^{(\mathcal{C})}/\mathrm{R}\) is called the Grothendieck group of C and is denoted by \(\mathrm{K}(\mathcal{C})\) . For a module E of type C, we denote the image of \(\mathrm{cl}(\mathrm{E})\) in \(\mathrm{K}(\mathcal{C})\) by \([\mathrm{E}]_{\mathcal{C}}\) (or sometimes {[}E{]} when there is ambiguity about C). We then have the following universal property.

PROPOSITION 4. --- a) The mapping \(E \mapsto [E]_{\mathcal{C}}\) from C to \(\mathrm{K}(\mathcal{C})\) is additive.

\begin{enumerate}
\def\labelenumi{\alph{enumi})}
\setcounter{enumi}{1}
\tightlist
\item
  Let G be an abelian group, and let \(\varphi : \mathcal{C} \to G\) be an additive function of modules. There exists a unique homomorphism \(u: K(\mathcal{C}) \to G\) such that we have \(\varphi(E) = u([E]_{\mathcal{C})}\) for every module E of type \(\mathcal{C}\) .
\end{enumerate}

Assertion a) is obvious. Let us prove b). There exists a homomorphism \(u^{\prime}: \mathbf{Z}^{(\mathcal{C})} \to \mathrm{G}\) that extends \(\varphi\) . Since \(\varphi\) is additive, we have \(u^{\prime}(r_{\mathcal{E}}) = 0\) for every exact sequence ( \(E\) ) of modules of type C; therefore, R is contained in the kernel of \(u^{\prime}\) . Consequently, when passing to the quotient, \(u^{\prime}\) defines a homomorphism u from \(\mathrm{K}(\mathcal{C})\) to G, and it is clear that we have \(\varphi(\mathrm{E}) = u([\mathrm{E}]_{\mathcal{C}})\) for every A-module of type C. The group \(\mathrm{K}(\mathcal{C})\) is generated by the set of elements \([E]_{C}\) for E running through C; the uniqueness of u follows.

Let C and D be additive sets of classes of A-modules such that \(C \subset D\) . Since the mapping \(E \mapsto [E]_{D}\) from nC to the Grothendieck groupnK(D) is additive, there exists a homomorphism \(\gamma_{\mathcal{D},\mathcal{C}} : K(\mathcal{C}) \to K(\mathcal{D})\) , called canonical, characterized by the formula \(\gamma_{\mathcal{D},\mathcal{C}}([E]_{\mathcal{C})} = [E]_{\mathcal{D}}\) for every module E of type C. It is not always injective (VIII, p.~210, Exercise 13).

Example. --- *Let A be a Noetherian commutative ring, and let \(\Sigma\) be its spectrum. For any integer \(n \geqslant 0\) , denote by \(\mathcal{C}^{\geqslant n}\) the set of classes of finitely generated A-modules with support of codimension \(\geqslant n\) in \(\Sigma\) . Set \(\mathrm{K}(n,\mathrm{A}) = \mathrm{K}(\mathcal{C}^{\geqslant n})\) and \(\gamma_{n} = \gamma_{\mathcal{C}^{\geqslant n},\mathcal{C}^{\geqslant n + 1}}\) . We have a sequence of homomorphisms

\[
\gamma_ {n}: \mathrm{K} (n + 1, \mathrm{A}) \longrightarrow \mathrm{K} (n, \mathrm{A}).
\]

We can prove (AC, VIII, §1, n° 5, p.~13, proposition 10) that in K(n, A), the elements \([A/p]_{C_{n}}\) , where p runs through the set of elements of \(\Sigma\) of height n, form a basis of a Z-module supplementary to the image of \(\gamma_{n}\) . More precisely, for every module E of type \(C \geqslant n\) , we have

\[
[ \mathrm{E} ] _ {\mathcal {C} ^ {\geqslant n}} \equiv \sum_ {\{\mathfrak {p} \in \Sigma | \operatorname{ht} (\mathfrak {p}) = n \}} \operatorname{long} _ {\mathrm{A} _ {\mathfrak {p}}} (\mathrm{E} _ {\mathfrak {p}}) \cdot [ \mathrm{A} / \mathfrak {p} ] _ {\mathcal {C} ^ {\geqslant n}} \quad (\mathrm{mod} \operatorname{Im} \gamma_ {n}). _ {*}
\]

The group \(\mathrm{K}(\mathcal{C})\) is generated by the elements of the form \([E]_{\mathcal{C}}\) with \(E \in C\) , and we have \([E \oplus E']_{\mathcal{C}} = [E]_{\mathcal{C}} + [E']_{\mathcal{C}}\) by relation (1) (VIII, p.~184). Every element of \(\mathrm{K}(\mathcal{C})\) is therefore of the form \([E]_{\mathcal{C}} - [F]_{\mathcal{C}}\) , where E and F belong to C.

An element of \(\mathrm{K}(\mathcal{C})\) is called effective if it is of the form \([E]_{\mathcal{C}}\) for an A-module E of type C. The set of effective elements of \(\mathrm{K}(\mathcal{C})\) is denoted by \(\mathrm{K}(\mathcal{C})^{+}\) ; it is a submonoid of \(\mathrm{K}(\mathcal{C})\) , and \(\mathrm{K}(\mathcal{C})\) can be identified with the group of differences of \(\mathrm{K}(\mathcal{C})^{+}\) (I, §2, No.~4, p.~20).

PROPOSITION 5. --- Let E and F be modules of type C. The equality \([E]_{C} = [F]_{C}\) holds if and only if there exist exact sequences of modules of type C

(ε)

\[
0 \longrightarrow \mathrm{L} \longrightarrow \mathrm{P} \longrightarrow \mathrm{M} \longrightarrow 0,\tag{\( \mathcal{F} \)}
\]

\[
0 \longrightarrow \mathrm{L} \longrightarrow \mathrm{Q} \longrightarrow \mathrm{M} \longrightarrow 0
\]

such that \(\mathrm{E} \oplus \mathrm{Q}\) is isomorphic to \(\mathrm{F} \oplus \mathrm{P}\) .

The stated condition is sufficient because it implies the relations

\[
[ \mathrm{P} ] _ {\mathcal {C}} = [ \mathrm{L} ] _ {\mathcal {C}} + [ \mathrm{M} ] _ {\mathcal {C}}, \qquad [ \mathrm{Q} ] _ {\mathcal {C}} = [ \mathrm{L} ] _ {\mathcal {C}} + [ \mathrm{M} ] _ {\mathcal {C}}, \qquad [ \mathrm{E} ] _ {\mathcal {C}} + [ \mathrm{Q} ] _ {\mathcal {C}} = [ \mathrm{F} ] _ {\mathcal {C}} + [ \mathrm{P} ] _ {\mathcal {C}},
\]

which give \([\mathrm{E}]_{\mathcal{C}} = [\mathrm{F}]_{\mathcal{C}}\)

Now suppose that we have \([E]_{C} = [F]_{C}\) . By the construction of the group \(\mathrm{K}(\mathcal{C})\) , there exist two finite families of exact sequences of modules of type C

\[
0 \longrightarrow \mathrm{G} _ {i} ^ {\prime} \longrightarrow \mathrm{G} _ {i} \longrightarrow \mathrm{G} _ {i} ^ {\prime \prime} \longrightarrow 0\tag{\((\mathcal{G}_i)\)}
\]

for \(i\in \mathrm{I}\) and

\[
0 \longrightarrow \mathrm{H} _ {j} ^ {\prime} \longrightarrow \mathrm{H} _ {j} \longrightarrow \mathrm{H} _ {j} ^ {\prime \prime} \longrightarrow 0\tag{\((\mathcal{H}_j)\)}
\]

for \(j\in \mathrm{J}\) such that in \(\mathbf{Z}^{(\mathcal{C})}\) , we have

\[
\operatorname{cl} (\mathrm{E}) - \operatorname{cl} (\mathrm{F}) = \sum_ {j \in \mathrm{J}} r _ {\mathcal {H} _ {j}} - \sum_ {i \in \mathrm{I}} r _ {\mathcal {G} _ {i}}.
\]

More explicitly, this relation can be written as

\[
\begin{array}{c} \operatorname{cl} (\mathrm{E}) + \sum_ {i \in \mathrm{I}} \operatorname{cl} (\mathrm{G} _ {i}) + \sum_ {j \in \mathrm{J}} \operatorname{cl} (\mathrm{H} _ {j} ^ {\prime}) + \sum_ {j \in \mathrm{J}} \operatorname{cl} (\mathrm{H} _ {j} ^ {\prime \prime}) \\ = \operatorname{cl} (\mathrm{F}) + \sum_ {i \in \mathrm{I}} \operatorname{cl} (\mathrm{G} _ {i} ^ {\prime}) + \sum_ {i \in \mathrm{I}} \operatorname{cl} (\mathrm{G} _ {i} ^ {\prime \prime}) + \sum_ {j \in \mathrm{J}} \operatorname{cl} (\mathrm{H} _ {j}). \end{array}\tag{7}
\]

Set \(G = \bigoplus_{i \in I} G_i\) , \(G' = \bigoplus_{i \in I} G_i'\) , etc. By passing to the direct sums, we obtain the exact sequences

\begin{enumerate}
\def\labelenumi{(\Alph{enumi})}
\setcounter{enumi}{6}
\tightlist
\item
\end{enumerate}

\[
0 \longrightarrow \mathrm{G} ^ {\prime} \stackrel {p} {\longrightarrow} \mathrm{G} \stackrel {q} {\longrightarrow} \mathrm{G} ^ {\prime \prime} \longrightarrow 0,\tag{H}
\]

\[
0 \longrightarrow \mathrm{H} ^ {\prime} \stackrel {r} {\longrightarrow} \mathrm{H} \stackrel {s} {\longrightarrow} \mathrm{H} ^ {\prime \prime} \longrightarrow 0
\]

consisting of modules of type \(\mathcal{C}\) .

Next, let \(M_{1},\ldots,M_{m},N_{1},\ldots,N_{n}\) be A-modules of type C. If we have \(\sum_{i=1}^{m}\operatorname{cl}(M_{i})=\sum_{j=1}^{n}\operatorname{cl}(N_{j})\) in the group \(\mathbf{Z}^{(\mathcal{C})}\) , then we have m=n and there exists a permutation \(\sigma\in\mathfrak{S}_{m}\) such that \(\operatorname{cl}(M_{i})=\operatorname{cl}(N_{\sigma(i)})\) for every \(1\leqslant i\leqslant m\) (I, §7, No.~9, p.~95, Proposition 11). Consequently, the modules \(\bigoplus_{i=1}^{m}M_{i}\) and \(\bigoplus_{j=1}^{n}N_{j}\) are isomorphic. In particular, from (7), we deduce the existence of an isomorphism from \(E\oplus Q\) to \(F\oplus P\) , where we have set

\[
\mathrm{P} = \mathrm{G} ^ {\prime} \oplus \mathrm{G} ^ {\prime \prime} \oplus \mathrm{H}, \quad \mathrm{Q} = \mathrm{G} \oplus \mathrm{H} ^ {\prime} \oplus \mathrm{H} ^ {\prime \prime}.
\]

We also set

\[
\mathrm{L} = \mathrm{G} ^ {\prime} \oplus \mathrm{H} ^ {\prime}, \quad \mathrm{M} = \mathrm{G} ^ {\prime \prime} \oplus \mathrm{H} ^ {\prime \prime}.
\]

The modules L, M, P, Q are of type C, and the sequence

\[
0 \longrightarrow \mathrm{L} \stackrel {\lambda} {\longrightarrow} \mathrm{P} \stackrel {\mu} {\longrightarrow} \mathrm{M} \longrightarrow 0,\tag{\( \mathcal{E} \)}
\]

where we define \(\lambda\) and \(\mu\) by

\[
\lambda (g ^ {\prime}, h ^ {\prime}) = \left(g ^ {\prime}, 0, r (h ^ {\prime})\right), \quad \mu (g ^ {\prime}, g ^ {\prime \prime}, h) = \left(g ^ {\prime \prime}, s (h)\right),
\]

is exact. We construct an exact sequence

\[
0 \longrightarrow \mathrm{L} \longrightarrow \mathrm{Q} \longrightarrow \mathrm{M} \longrightarrow 0\tag{\( \mathcal{F} \)}
\]

in the same way; this concludes the proof.

The set C is a commutative monoid for the law of composition \((\mathrm{E}, \mathrm{E}^{\prime}) \mapsto \mathrm{cl}(\mathrm{E} \oplus \mathrm{E}^{\prime})\) . We sometimes denote the group of differences of the commutative monoid \(\mathcal{C}\) (I, §2, No.~4, p.~20) by \(\mathrm{K}'(\mathcal{C})\) and call it the Grothendieck group of \(\mathcal{C}\) for direct sums. For any module E of type \(\mathcal{C}\) , we denote the image of \(\operatorname{cl}(\mathrm{E})\) in \(\mathrm{K}'(\mathcal{C})\) by \([\mathrm{E}]_{\mathcal{C}}'\) .

PROPOSITION 6. --- a) The mapping \(E \mapsto [E]_{\mathcal{C}}'\) from C to \(K'(\mathcal{C})\) is a weakly additive function of modules.

\begin{enumerate}
\def\labelenumi{\alph{enumi})}
\setcounter{enumi}{1}
\item
  Let G be an abelian group, and let \(\varphi : \mathcal{C} \to G\) be a weakly additive function of modules. There exists a unique group homomorphism \(u: K'(\mathcal{C}) \to G\) such that we have \(\varphi(E) = u([E]_{\mathcal{C}}')\) for every module E of type \(\mathcal{C}\) .
\item
  Let \(\mathrm{E}\) and \(\mathrm{F}\) be modules of type \(\mathcal{C}\) . We have \([\mathrm{E}]_{\mathcal{C}}' = [\mathrm{F}]_{\mathcal{C}}'\) if and only if there exists a module \(\mathrm{M}\) of type \(\mathcal{C}\) such that \(\mathrm{E} \oplus \mathrm{M}\) is isomorphic to \(\mathrm{F} \oplus \mathrm{M}\) .
\end{enumerate}

Assertion a) is obvious. Assertion b) follows from (I, §2, No.~4, p.~19, Theorem 1) and assertion c) from (I, §2, No.~4, p.~18, Proposition 6).

Since the mapping \(E \mapsto [E]_{\mathcal{C}}\) from C to \(\mathrm{K}(\mathcal{C})\) is a weakly additive function of modules, we can deduce a homomorphism \(u : \mathrm{K}'(\mathcal{C}) \to \mathrm{K}(\mathcal{C})\) from it. This homomorphism is surjective but is not always an isomorphism (VIII, p.~191, Remark 2).

Let \(R'\) be the subgroup of \(\mathbf{Z}^{(\mathcal{C})}\) generated by the elements of the form \(r_{\mathcal{E}}\) , where \(\mathcal{E}\) is a split exact sequence of A-modules of type \(\mathcal{C}\) , that is, by the elements of the form \(\operatorname{cl}(\operatorname{E}' \oplus \operatorname{E}'') - \operatorname{cl}(\operatorname{E}') - \operatorname{cl}(\operatorname{E}'')\) , where \(E'\) and \(E''\) are modules of type \(\mathcal{C}\) . The canonical mapping from \(\mathcal{C}\) to the quotient group \(\mathbf{Z}^{(\mathcal{C})}/R'\) extends to a group homomorphism \(v: K'(\mathcal{C}) \to \mathbf{Z}^{(\mathcal{C})}/R'\) . This is an isomorphism. Indeed, the canonical mapping from \(\mathcal{C}\) to \(K'(\mathcal{C})\) extends to a group homomorphism from \(\mathbf{Z}^{(\mathcal{C})}\) to \(K'(\mathcal{C})\) whose kernel contains \(R'\) , and therefore to a homomorphism \(v': \mathbf{Z}^{(\mathcal{C})}/R' \to K'(\mathcal{C})\) by passing to the quotient; it is clear that \(v\) and \(v'\) are inverse bijections.

An element of \(\mathrm{K}^{\prime}(\mathcal{C})\) is called effective if it is of the form \([E]_{\mathcal{C}}^{\prime}\) for an A-module E of type C. The set of effective elements of \(\mathrm{K}^{\prime}(\mathcal{C})\) is denoted by \(\mathrm{K}^{\prime}(\mathcal{C})^{+}\) .

